\subsection{A canonical higher exterior-power identity}
\label{sec:higher-car}

Higher creation and annihilation operators suggest a low-dimensional
intermediate exterior power. Their partial contractions must also be
charged. The following exact calculation groups all those corrections
into one opposite-cut operator and then bounds its actual rank.
It concerns this canonical grouping only.

Work over a field $K$ of characteristic two. For $L\ge2p\ge2$, put
$U=\Lambda^pK^L$ and $Q=\binom Lp$, using the standard subset bases.
Let $\mu:U\otimes U\to\Lambda^{2p}K^L$ be wedge multiplication and put
$B=\mu^{\mathsf T}\mu$. Let $a_S$ remove the subset $S$ when present
and give zero otherwise. Define
\[
\begin{gathered}
 D_r=\sum_{|S|=r}a_S\otimes a_S:
 U\otimes U\longrightarrow
 \Lambda^{p-r}K^L\otimes\Lambda^{p-r}K^L,\\
 K_r=D_r^{\mathsf T}D_r,
 \qquad N=\sum_{r=1}^pK_r.
\end{gathered}
\]
Transposes here use the displayed bases, and $\Gamma$ is ordinary
partial transpose on the second $U$.

\begin{proposition}[Exact canonical contraction accounting]
\label{prop:higher-car-grouping}
One has
\[
 I=B+\Gamma N,\qquad \operatorname{rank}B=\binom L{2p}.
\]
Let $\mathrm{KG}(n,b)$ be the matrix on $b$-subsets of $[n]$ whose
entry is one precisely for disjoint pairs. Then
\begin{equation}\label{eq:higher-car-rank}
 \operatorname{rank}N
 =\sum_{a=0}^{p-1}\binom L{2a}\binom{2a}a\,
   \operatorname{rank}\bigl(I+\mathrm{KG}(L-2a,p-a)\bigr).
\end{equation}
These identities hold over every characteristic-two extension field.
Along every sequence with $L\ge2p\ge2$ and $Q\to\infty$,
\begin{equation}\label{eq:higher-car-exponent}
 \liminf\frac{\log(\operatorname{rank}B+\operatorname{rank}N)}
                  {\log Q}
 \ge\log_2 3>\sqrt2.
\end{equation}
\end{proposition}
\begin{proof}
A basis pair $(I,J)$ has symmetric difference $A=I\mathbin\triangle J$, of size $2a$,
a fixed orientation $I|_A$ of size $a$, and common subset
$C=I\cap J$ of size $b=p-a$. Each $D_r$ preserves the symmetric difference and its
orientation. On the remaining $n=L-2a$ coordinates it is the inclusion
matrix from $b$-subsets to $(b-r)$-subsets. Thus the corresponding
$K_r$ block has entry
$\binom{|C\cap C'|}{b-r}$. If $b\ge1$, summing over $1\le r\le b$
gives one on the diagonal and on disjoint pairs, and zero otherwise:
for $t=|C\cap C'|<b$ the sum is $2^t$, while for $t=b$ it is $2^b-1$.
The block is therefore $I+\mathrm{KG}(n,b)$; when $b=0$ it is zero.
There are $\binom L{2a}\binom{2a}a$ choices of symmetric difference and orientation,
proving~\eqref{eq:higher-car-rank}.

Include $K_0=I$. An entry of $\Gamma\sum_{r=0}^pK_r$ with row
$(I,J)$ and column $(K,T)$ vanishes unless
$I+J=K+T$ as integer occupancy vectors. Under that condition the
intersection of the common subsets in the transposed block is
$C_0=I\cap J=K\cap T$. Summing all the allowed $r$ counts every
subset of $C_0$, so the entry is $2^{|C_0|}$. It is therefore one
exactly when the two pairs are disjoint and have the same union.
This is $B$. Hence $B=\Gamma(I+N)$, which gives the claimed identity.
The matrix $B$ consists of nonempty all-one blocks indexed by the
$2p$-element union, so its rank is $\binom L{2p}$.

Every block $I+\mathrm{KG}(n,b)$ with $b\ge1$ has diagonal entries
one, and consequently
\[
 \operatorname{rank}N\ge
 \sum_{a=0}^{p-1}\binom L{2a}\binom{2a}a.
\]
Write $h=\log_2 3$, let $H$ denote base-two Shannon entropy with
$H(x)=H(x,1-x)$, and put
$T(x)=H(2x)+2x=H(x,x,1-2x)$.
The multinomial type bound gives, uniformly in $a$,
\[
 \log_2\left(\binom L{2a}\binom{2a}a\right)
 \ge L T(a/L)-O(\log L).
\]
For $0\le x\le1/3$, one has $T(x)\ge hH(x)$.
Indeed $T-hH$ is concave, since its second derivative on the
interior is
\[
 \frac{(h-2)+(2-2h)x}
      {(\ln2)x(1-x)(1-2x)}<0,
\]
and its endpoint values are zero and $h(1-H(1/3))>0$.
If $p/L\le1/3$, choose $a=p-1$; changing the entropy argument by
$1/L$ costs $O(\log L)$ after multiplication by $L$.
If $p/L\ge1/3$, choose $a=\min(p-1,\lfloor L/3\rfloor)$ and use
$T(1/3)=h$ and $Q\le2^L$. In both cases,
\[
 \log_2\operatorname{rank}N\ge h\log_2Q-O(\log L).
\]
When $p\to\infty$, the error is negligible because
$\log Q/\log L\to\infty$: use $Q\ge(L/p)^p$ and split at
$p=\sqrt L$. On a subsequence with bounded $p$, pass further to a
constant $p$; then $\log\binom L{2p}/\log\binom Lp\to2$.
These two cases prove the liminf assertion.
Finally $3^2>2^3$ gives $h>3/2>\sqrt2$.
\end{proof}

No eigenvalue argument or modular diagonalizability is used in the
rank lower bound. Other coefficients of the $K_r$, products of
contraction operators, and unrestricted crossing decompositions
remain outside this calculation.
